\section{计算思维}

\begin{frame}{什么是计算思维？}
\begin{itemize}                                                   
  \item {\cbf{推动人类文明进步和科技发展的三大科学支柱}}
    \begin{itemize}
      \item 理论科学
      \item 实验科学
      \item 计算科学
    \end{itemize}
  \item \cbf{人类认识世界和改造世界的三种思维}
    \begin{itemize}
      \item 逻辑思维
      
      以推理和演绎为特征，以数学学科为代表。
      
      \item 实证思维
      
      以观察和总结自然规律为特征，以物理等学科为代表。
      \item 计算思维
      
      以设计和构造为特征，以计算机学科为代表。
    \end{itemize}
\end{itemize}
\end{frame}


\begin{frame}
    \frametitle{什么是计算思维？}
    \begin{minipage}{0.68\linewidth}
        计算思维是运用计算机科学的基础概念去求解问题、设计系统和理解人类行为，它涵盖了计算机科学的一系列思维活动。
        \begin{itemize}
          \item 问题求解中的计算思维
          \item 系统设计中的计算思维
          \item 人类行为理解中的计算思维
        \end{itemize}

计算思维如同所有人都具备“读、写、算”能力一样，成为适合于每个人的一种普遍的认识和一类普适的技能。

    \end{minipage}\hfill
    \begin{minipage}{0.28\linewidth}
        \begin{figure}[htpb]
            \centering
            \includegraphics[width=\linewidth]{周以真}
                \caption{周以真\\卡内基梅隆大学教授}
        \end{figure}
    \end{minipage}
\end{frame}


\begin{frame}
    \frametitle{理解计算思维，首先要理解计算}
    理解计算思维的前提是理解计算，研究在解决问题过程中，哪些是可计算的，以及如何计算。

    \begin{itemize}
      \item 算术运算
      \item 集合运算
      \item 逻辑运算
      \item 条件运算
      \item $\cdots \cdots$
    \end{itemize}

    当我们跳出算术运算的局限，理解了计算的本质后，就会发现原来好多看似不可计算的东西都能变得可计算，也就很容易理解计算思维的普适性了。

\end{frame}


\begin{frame}
    \frametitle{计算思维的本质}
    \begin{minipage}{0.48\linewidth}
        \begin{itemize}
            \item \bf{抽象}
            
            把实际问题抽象为数学问题，并建模
              \begin{itemize}
                \item 问题建模：忽略某些细节，如最短路径求解，方程计算。
                \item 数据建模：考虑数据特点，数据可抽象为线性结构，树结构，图结构等。
              \end{itemize}
            \item \bf{自动化}
            
            机械地一步一步地自动执行，并选择合适的计算机执行问题的抽象。
          \end{itemize}
    \end{minipage}\hfill
    \begin{minipage}{0.48\linewidth}
        \begin{figure}[htpb]
            \centering
            \includegraphics[width=\linewidth]{七桥问题}
            \caption{哥尼斯堡的七座桥}
        \end{figure}
        \begin{example}
            大象放冰箱里分几步？
        \end{example}
    \end{minipage}
    
\end{frame}


\begin{frame}
    \frametitle{问题求解中的计算思维}
    \begin{figure}[htpb]
        \centering
        \includegraphics[width=.8\linewidth]{问题求解}
        \caption{使用计算机进行问题求解的计算思维过程}
    \end{figure}
\end{frame}


\begin{frame}
    \frametitle{系统设计中的计算思维}
    \begin{figure}[htpb]
        \centering
        \includegraphics[width=\linewidth]{系统设计}
        \caption{将系统中的变换过程映射为符号变换---形式化，然后用计算机进行处理。}
    \end{figure}
\end{frame}


\begin{frame}
    \frametitle{系统设计中的计算思维}
    我们经常在较低抽象级之上再建立一个较高的抽象级，以便隐藏低抽象级的复杂 细节，提供更简单的求解方法。
        \begin{blankblock}
            例如，对计算本身的理解就可以形成“电子电路$\to$门逻辑$\to$二进制$\to$机器语言指令$\to$高级语言程序”这样一个由低到高的抽象层次，我们之所以在高级 语言程序这个层次上学习计算，当然是为了隐藏那些低抽象级的繁琐细节。
        \end{blankblock}
        \begin{blankblock}{}
            又如，在互联网 上发送一封电子邮件实际上要经过不同抽象级的多层网络协议才得以实现，写邮件的人肯定 不希望先掌握网络低层知识才能发送邮件。
        \end{blankblock}
\end{frame}


\begin{frame}
    \frametitle{人类行为理解中的计算思维}
    计算思维不仅能够反映人类思维活动，还能够通过计算手段来研究人类的行为。
    
    利用计算手段来研究人类的行为，可视为社会计算，即通过各种信息技术手段，设计、实施和评估人与环境之间的交互。

    \begin{itemize}
      \item 网购习惯
      \item 推荐算法
    \end{itemize}

\end{frame}


\begin{frame}
    \frametitle{计算思维理论的发展}
\begin{table}[]
    \footnotesize
    \begin{tabularx}{\textwidth} {
         >{\raggedright\arraybackslash}m{0.4\textwidth}
         >{\raggedright\arraybackslash}m{0.5\textwidth}
         }
    \toprule
    计算思维1.0                              & 计算思维2.0                                            \\ \midrule
    计算思维是计算机科学家的思维模式                     & 计算思维是所有领域的科学家在应用计算和计算模型时所采取的思维模式                   \\ \hline
    计算思维是从计算机科学中产生 ， 逐步扩散到其他科学领域         & 每个科学领域都有属于自身特点的计算思维，贯穿通过计算模型求解的设计 、 执行和评价过程        \\ \hline
    学习编程能力过程中产生计算思维 ， 计算思维主要是大量的编程技巧     & 计算思维提供了有关模型和算法的丰富概念 ， 学习和理解计算思维用以产生科学的编程能力         \\ \hline
    算法必须与程序紧密联系                          & 算法并不一定要联系程序和计算机 ， 算法是信息处理的技术 ， 一些生物体内部进行的自然计算也属于算法 \\ \hline
    实施算法设计要求具备计算机科学的领域知识 ， 特别是关于计算机的基本知识 & 每个人都可以在本领域找到问题计算求解的模型 、 算法和操作 ， 并不需要图灵机这样的通用概念     \\ \hline
    应该关心程序背后的计算机概念而不是程序本身                & 应该关心程序背后反映领域特点的内容 ， 注重本学科解决问题的计算化趋势                \\ \bottomrule
    \end{tabularx}
    \end{table}
\end{frame}


\begin{frame}
    \frametitle{计算思维的特征}
    \begin{itemize}
      \item 概念化，不是程序化
      \item 基础的，不是机械的技能
      \item 人的，不是计算机的思维
      \item 数学和工程思维的互补与融合
      \item 是思想，不是人造品
      \item 不能超出计算机的能力范围
    \end{itemize}
\end{frame}


\begin{frame}[fragile]
    \frametitle{不能超出计算机的能力范围}
    计算机的能力：能够高速执行大量指令

    计算机的限制：每条指令只能进行有限的一些简单操作

    \begin{example}
        某同学刚学了几天程序语言设计，想试试自己的能力，求方程$3x^2+5x-15=0$的根。
该同学写出下述的程序段:
    \end{example}
    \begin{codeblock}{}
    3*x*x + 5*x - 15 = 0；# 给出方程，让计算机求解
    print(x); # 输出结果
    \end{codeblock}
\end{frame}


\begin{frame}[fragile]
    \frametitle{}
    \begin{minipage}{0.48\linewidth}
        \begin{codeblock}{}
    a = float(input("输入a:"))   
    b = float(input("输入b:"))
    c = float(input("输入c:"))
    #计算delta
    delta = b**2-4*a*c
    #判断解的个数
    if delta < 0:
        print("方程无解！")
    elif delta == 0:
        x = b/(-2*a)
        print("x1=x2=",x)
    else:
        x1 = (b+delta**0.5)/(-2*a)
        x2 = (b-delta**0.5)/(-2*a)
    #输出x1,x2
    print("x1=",x1)
    print("x2=",x2)
            \end{codeblock}
    \end{minipage}\hfill
    \begin{minipage}{0.48\linewidth}
        \Large
        \[
            \text{方程} ax^2 + bx + c = 0 \text{的解为}
        \]

        \[
            \frac{-b \pm \sqrt{b^2 - 4ac}}{2a}
        \]
    \end{minipage}
    
\end{frame}


\begin{frame}
    \frametitle{计算思维案例：word邮件合并}
    \begin{minipage}{0.48\linewidth}
        \begin{figure}[htpb]
            \centering
            \includegraphics[width=\linewidth]{邮件合并}
        \end{figure}
    \end{minipage}\hfill
    \begin{minipage}{0.48\linewidth}
        工资条，准考证，成绩单，证书等一类问题

        \begin{enumerate}
          \item 抽象
          
          对共同的特征进行抽取，建立模型
            \begin{itemize}
              \item 文档固定模版：统一风格
              \item 文档内容分为固定内容和动态内容
              \item 动态内容具有相同数据结构
            \end{itemize}
          \item 自动化
          
          \begin{itemize}
            \item 做好数据源，文档模板
            \item 将数据每个属性写入到文档模板
            \item 重复步骤B，合并所有数据
          \end{itemize}
        \end{enumerate}
    \end{minipage}
\end{frame}


\begin{frame}
    \frametitle{计算思维对其他学科的影响}
    \begin{itemize}
      \item 数学 
      \item 化学 
      \item 生物学
      \item 社会学
      \item 博弈论
    \end{itemize}
\end{frame}


\begin{frame}
    \frametitle{计算数学}
    计算机对数学来说过去只是一个数值计算工具，用于快速、大规模的数值计算，对数值计算方法的研究导致了计算数学的形成。

后来数学家利用计算机进行代数演算，形成了计算机代数；利用计算机研究几何问题，形成了计算几何学。

数学家还利用计算机去验证数学猜想，虽然不能证明猜想，但是一旦发现反例就可以推翻猜想，以免数学家毕生投入到一个不成立的猜想之中。

总之，现在计算机已经成为数学的研究手段，大大扩展了数学家的能力。

\end{frame}


\begin{frame}
    \frametitle{机器证明}
    机器证明是使用计算机证明定理，也称为定理的机械证明或自动证明。

    70年代后期，在计算机技术大发展的背景下，吴文俊继承和发展了中国古代数学的传统（即算法化思想），转而研究几何定理的机器证明，彻底改变了这个领域的面貌，是国际自动推理界先驱性的工作，被称为“吴方法”，产生了巨大影响。
    \begin{figure}[htpb]
        \centering
        \includegraphics[width=.3\linewidth]{吴文俊}
        \caption{中科院资深院士，中国人工智能之父---吴文俊}
    \end{figure}
\end{frame}


\begin{frame}
    \frametitle{计算化学}
    \begin{minipage}{0.58\linewidth}
        \begin{itemize}
            \item 数值计算。对化学各专业的数学模型进行数值计算或方程求解
            \item 化学模拟，数值模拟、过程模拟和实验模拟
            \item 在有机分析中根据图谱数据库进行图谱检索等
            \item 利用原子计算去探索化学现象
            \item 用优化和搜索算法寻找优化化学反应条件和提高产量的物质
          \end{itemize}
    \end{minipage}\hfill
    \begin{minipage}{0.38\linewidth}
        \begin{figure}[htpb]
            \centering
            \includegraphics[width=\linewidth]{轨道}
        \end{figure}
    \end{minipage}
    
\end{frame}


\begin{frame}
    \frametitle{计算生物学}
    \begin{minipage}{0.7\linewidth}
        \begin{itemize}
            \item 霰弹枪算法大大降低了人类基因组测序的成本，提高了测序的速度；
            \item 利用绳结来模拟蛋白质结构，用计算过程来模拟蛋白质动力学，并且运用数据挖掘与聚类分析的方法进行蛋白质结构的预测；\link{https://zhuanlan.zhihu.com/p/334267131}
            \item 开发了生物数据处理分析方法和知识库，帮助人们从分子层次上认识生命的本质及其进化规律。
            \item 在医学领域，机器人手术、借助于计算机的分析诊断及可视化系统在临床中已经广泛应用。
          \end{itemize}
    \end{minipage}\hfill
    \begin{minipage}{0.25\linewidth}
        \begin{figure}[htpb]
            \centering
            \includegraphics[width=\linewidth]{生物}
        \end{figure}
    \end{minipage}
\end{frame}


\begin{frame}
    \frametitle{谷歌AlphaFold解决蛋白质折叠问题}
        蛋白质的基本组成单位是氨基酸。蛋白质会由所含氨基酸残基的亲水性、疏水性、带正电、带负电等特性通过残基间的相互作用而折叠成一个立体的三级结构。

        这个折叠后的三维形状至关重要，因为它决定了蛋白质的工作方式。如果开发新药物的科学家知道蛋白质的形状，他就更容易找到一个可以与之结合的分子，使这个蛋白质可以改变它的行为。问题在于，很难预测蛋白质将会呈现哪种形状。

        1972年的诺贝尔化学奖得主安芬森提出的著名假设，蛋白质分子的一级结构决定其立体结构。也就是说，如果从蛋白质的一级结构就能知道它的三级立体结构，那么就可以直接从基因推测其编码蛋白质所对应的生物学功能。
    
        AlphaFold系统已经能够以惊人的速度和准确性预测出氨基酸串会折叠成什么样子的立体形状。
\end{frame}



\begin{frame}
    \frametitle{计算博弈论}

    博弈论是一种数学工具，它允许我们对特定环境下的利己主义和理性行为者的战略互动进行推理。该结构提供了一组框架，描述了在这样一个战略代理人的环境下的理性结果。虽然博弈论领域起源于经济文献，但计算机科学家在过去几十年里从建模和计算的角度对这一领域做出了重大贡献，这导致了计算博弈论(Algorithmic game theory)。

    和传统的博弈论和计算机科学相比，算法博弈论主要关注点在互联网网络，非传统拍卖和人类行为的根本问题

    \begin{itemize}
      \item 社交网络里的个体行为
      \item 用拍卖的方式出售关键字广告位
    \end{itemize}
\end{frame}


\begin{frame}
    \frametitle{计算社会学}
    

    计算社会学是当代社会学界借助计算机 、互联网与人工智能技术等现代科技手段，利用大数据 、新方法来获取数据与分析数据，从而研究与解释社会的一种新的范式或思维方式 。计算社会学使用密集演算的方法来分析与模拟社会现象。

根据社会科学理论，以计算技术为工具，可以研究人类社会的组成、关系、结构、层次、行为、运动等问题；研究社会媒体，促进社会人群的交流与合作。

\begin{itemize}
  \item 如何通过数据预测美国总统大选的结果？
  \item 流行病的发展轨迹
\end{itemize}

\end{frame}








